\section{Results}
\label{sec:results}

\subsection{Polsby-Popper Compactness}
\label{sec:results:pp}

Table~\ref{tab:pp-results} reports mean Polsby-Popper scores by resolution
tier. The key finding is that both resolution tiers produce mean PP exceeding
the congressional baseline of $\overline{PP} = 0.361$ from
\citet{dellaLibera2026mec}.

\begin{table}[h]
\centering\small
\begin{tabular}{lrrrr}
\toprule
Resolution tier & States & $\overline{PP}$ & $\min PP$ & Frac.\ $PP < 0.20$ \\
\midrule
Tract (38 states)       & 38 & 0.381 & 0.092 & 4.2\% \\
Block-group (12 states) & 12 & 0.401 & 0.118 & 2.7\% \\
All 50 states           & 50 & 0.386 & 0.092 & 3.8\% \\
\midrule
Congressional (B track) & 50 & 0.361 & 0.071 & 6.9\% \\
\bottomrule
\end{tabular}
\caption{Polsby-Popper compactness by resolution tier. All state-house tiers
  exceed the congressional baseline. Block-group resolution achieves higher
  compactness than tract resolution, consistent with the finer boundary
  definition available at block-group resolution.}
\label{tab:pp-results}
\end{table}

The higher compactness at block-group resolution reflects two effects:
(1) block-group boundaries respect smaller-scale geographic features such as
arterial roads and natural features that are not visible at tract resolution;
and (2) the larger number of units provides richer boundary optimisation.
Companion paper F.5 \citep{dellaLibera2026compactScale} provides a
decomposition of the PP advantage into scale and shape components.

\subsection{Population Balance}
\label{sec:results:balance}

All 50 state house maps achieve maximum population deviation
$\max_i |\delta_i| \leq 0.5\%$. The median maximum deviation across all
50 states is $0.31\%$. No state required more than 15 boundary-swap
iterations to achieve $\pm 0.5\%$ balance after the initial METIS partition.

Table~\ref{tab:pop-balance} reports population balance statistics for
selected states. The bloc-group-resolution states generally achieve
tighter balance (median $0.24\%$) than tract-resolution states (median $0.35\%$),
because the finer geographic units allow more precise boundary placement.

\begin{table}[h]
\centering\small
\begin{tabular}{llrrrr}
\toprule
State & Resolution & $k$ & $T^*$ (persons) & Max $|\delta|$ & Median $|\delta|$ \\
\midrule
New Hampshire & block\_group & 400 & 3{,}342   & 0.48\% & 0.19\% \\
Wyoming       & block\_group & 60  & 9{,}494   & 0.42\% & 0.21\% \\
Washington    & block\_group & 98  & 74{,}043  & 0.31\% & 0.14\% \\
Nebraska      & block\_group & 49  & 38{,}382  & 0.38\% & 0.17\% \\
Texas         & tract        & 150 & 188{,}779 & 0.44\% & 0.22\% \\
California    & tract        & 80  & 494{,}762 & 0.39\% & 0.18\% \\
\bottomrule
\end{tabular}
\caption{Population balance statistics for selected states. All states
  achieve $\leq 0.5\%$ maximum deviation, well within the federal
  \emph{Wesberry} standard and far within the \emph{Reynolds} standard
  of $\leq 5\%$ for state legislative maps.}
\label{tab:pop-balance}
\end{table}

\subsection{County Splits}
\label{sec:results:splits}

For the 38 states with publicly available enacted 2020 state house maps,
the algorithmic maps produce fewer county splits in 30 states (79\%),
the same number in 4 states (11\%), and more in 4 states (11\%).
The mean reduction is 18\% (range: $-3\%$ to $-41\%$).

\begin{table}[h]
\centering\small
\begin{tabular}{llrrr}
\toprule
State & $k$ & Enacted splits & Algorithm splits & Reduction \\
\midrule
Wisconsin   & 99  & 28  & 18  & 36\% \\
Ohio        & 99  & 35  & 22  & 37\% \\
Pennsylvania & 203 & 51  & 30  & 41\% \\
Virginia    & 100 & 19  & 15  & 21\% \\
Washington  & 98  & 14  & 11  & 21\% \\
Texas       & 150 & 62  & 48  & 23\% \\
California  & 80  & 11  & 10  & 9\%  \\
\bottomrule
\end{tabular}
\caption{County splits: enacted vs.\ algorithmic maps for selected states.
  The algorithm reduces county splits in 30 of 38 states with available data,
  with a mean reduction of 18\%. Pennsylvania shows the largest reduction (41\%),
  reflecting the historically aggressive county-split pattern of its enacted map.}
\label{tab:county-splits}
\end{table}

\subsection{Runtime}
\label{sec:results:runtime}

Table~\ref{tab:runtime-full} reports wall-clock runtimes for all 50 states.
Runtimes range from 8 seconds (Delaware, $k=41$, tract) to 420 seconds
(New Hampshire, $k=400$, block\_group). The median runtime is 67 seconds.
Total runtime for the full 50-state sweep is approximately 68 minutes on a
single CPU core. With 4 parallel workers (via \texttt{redist states --workers 4}),
this reduces to approximately 18 minutes.

Block-group resolution states are 3.2$\times$ slower on average than
tract-resolution states at comparable $k$ values, consistent with the
$\sim\!3\times$ larger adjacency graph at block-group resolution.
